% title: Sum of Squares of Binomial Coefficients
% short-title: Binomial Squares
% abstract: We record a formal proof in Lean 4 of the classical identity asserting that the sum of the squares of the binomial coefficients in the $n$th row of Pascal's triangle is the central binomial coefficient $\binom{2n}{n}$.
% subjclass: 05A10
% keywords: binomial coefficient, Vandermonde identity, Pascal triangle, Lean 4
% author: Mario Hernández M.

\subsection{Sum of Squares of Binomial Coefficients}

The central identity in this note is the Lean theorem
\lean{Biblioteca.Demonstrations.sum_of_squares_of_binomial_coefficients}. It
states that for every natural number $n$,
\[
  \sum_{i=0}^{n} \binom{n}{i}^{2} = \binom{2n}{n}.
\]

\begin{theorem}
For every $n \in \mathbb{N}$,
\[
  \sum_{i=0}^{n} \binom{n}{i}^{2} = \binom{2n}{n}.
\]
\end{theorem}

\begin{proof}
This is a standard specialization of Vandermonde's identity. In the symmetric
case $m=n$, the convolution formula
\[
  \binom{m+n}{k}
  =
  \sum_{i+j=k} \binom{m}{i}\binom{n}{j}
\]
becomes
\[
  \binom{2n}{n}
  =
  \sum_{i=0}^{n} \binom{n}{i}\binom{n}{n-i}.
\]
Using the symmetry $\binom{n}{n-i}=\binom{n}{i}$, one obtains the desired sum
of squares. The Lean proof delegates this argument to the library theorem
\texttt{Nat.sum\_range\_choose\_sq}, which is itself proved in the Vandermonde
development of \texttt{Mathlib}.
\end{proof}
